% !TeX root = ../papers/nc-p3.tex
\section{四頂点切断における点加群と分離三次曲面の特異点}
\label{sec:finite-cut-point-pairings}
\begin{theorem}\label{thm:finite-cut-point-pairings}
四辺が有限な標準有限Frobenius候補を取り、連続四頂点 $i,\ldots,i+3$ の経路代数を $B$ とする。
点グラフの点 $p$ は、各頂点一次元、三辺の作用が $p_i,p_{i+1},p_{i+2}$ である右 $B$ 加群 $M(p)$ を定める。
元の無限代数のAS性によらず、$\nu_B M(p)\simeq M(p)[3]$ である。
異なる点 $p,q$ について
\[
 \dim\operatorname{Ext}^1_B(M(p),M(q))
 =\dim\operatorname{Ext}^2_B(M(p),M(q))
 =2-\operatorname{rank} w(p_{i+3},-,q_{i+1},-)
\]
であり、他の次数は零である。この階数は $p,q$ の交換にも切断の始点にも依存しない。
自己Extの次数0,3は一次元、次数1,2はその局所点の埋込み次元である。
\end{theorem}
\begin{proof}
$H_j=\ker p_j$、$P_j=e_jB$ とし、$h=w(p_{i+3},-,-,-)$ は添字を並べ直して
$H_i\otimes V_{i+1}\otimes H_{i+2}$ と読む。次は右加群の線型射影解消である。
\[
 0\to P_{i+3}\to H_{i+2}^*\otimes P_{i+2}
 \to H_i\otimes P_{i+1}\to P_i\to M(p)\to0 .
\]
両端の微分は指定した超平面の三文字で、中央は $h$ による。
第二自由項の係数は $R_i\cap(H_i\otimes V_{i+1})$ に等しい。
実際 $p_i$ での評価が階数3で、$h$ の末尾縮約がこの核内の独立な三列を与える。
最初の余核は点加群である。先頭の三文字を消した後は関係の評価が
次の点の消滅三平面全体を与え、残りの各段でも同じことを繰り返せる。
長さ2まではこの核の記述で完全である。長さ3での第二微分の核は
三次重なりのうち第一位置が $H_i$ に入る元である。全三次重なりは $w(s)$ であり、
条件は $w(s,p_i)=0$ なので $s$ は唯一の前点 $p_{i+3}$ の倍である。
長さ3の各項次元は左から $1,12,30,20,1$ であり、これで完全性が従う。
左加群についても同じ解消があり、右の表示に $\operatorname{Hom}_B(-,B)$ を施したものは
$h$ の中央行列を逆順に読んだその左解消である。従ってコホモロジーは次数3の左点加群だけで、
体上の双対を取ればSerre関手の式を得る。

右解消に $\operatorname{Hom}_B(-,M(q))$ を施すと次元 $1,3,3,1$ の列になり、
中央の微分が表示の反対切断である。異なる点なら両端は階数1で、同じ点なら両端は零である。
自己第一Extの次元は点グラフの接次元である。実際、双対数上への加群の持上げは各頂点で
自由階数1なので基底を選んで三つの非零な矢評価を持ち上げることと同じであり、
頂点基底変更の有効な三倍率は各評価の射影化と一致する。二辺の零点条件で三文字の列は
共通有限スキームと同型である。自己第二Extは同じ次元となる。
交換の不変性はSerre双対とこの二次数の一致による。

位相の不変性には相補切断の完全対合を使う。例えば位相0と2で、
$p_0,-$ と $q_2,-$ により得る二つの三空間の対合の階数が表示の階数である。
それぞれの直交補は辺 $p_3\to p_0, q_1\to q_2$ を使い $-,p_3$ と $-,q_1$ で得る三空間である。
六空間間の完全対合における三空間の対合と直交補の対合は同階数なので
位相3と1の反対切断も同じ階数である。巡回して結論を得る。
\end{proof}

\begin{corollary}\label{cor:finite-separator-singularity-graph}
全ての中央支持が4で二十点が被約なら、いずれの位置の二十点の配置からでも同じ無向グラフが得られる。
各点 $p$ に対し他の十九点で消える唯一の三次曲面を取り、
別の点 $q$ がその特異点となるとき $p,q$ を結ぶ。各頂点の次数は4以下で、
辺は切断上で点加群間の第一Extが一次元となることと同値である。
十五次元混合成分の一般点では辺はなく、二十の点加群は互いに全次数で直交する
三次元球面的対象である。さらにその直和は四頂点切断の忠実表現である。
\end{corollary}
\begin{proof}
分離三次曲面は定理\ref{thm:finite-cubic-separators}の行列式である。異点 $q$ における行列は
階数1または2である。0は全中央支持4に反する。階数1なら余因子は全て零である。
階数2なら、その三行三列での左右核は二つの隣り合う位相における
$q$ の座標を $p$ から射影したものである。従って余因子による微分は、倍率を除き
$w(p_0,q_1,-,q_3)$（曲面の位置を2とした）である。$q_1,q_3$ の反対切断は
$q$ が単純なので階数3で左核が $q_0$ のみである。$p_0\ne q_0$ なので微分は非零である。
これで前定理の交換・位相の不変性とExtの式がグラフに移る。

一つの四次元切断の階数一スキームが有限なら長さは4以下
（定理\ref{thm:finite-marked-secant-pencil}のSegreの完全交叉の計算）。
正次元なら一つの被約直線である（定理\ref{thm:finite-marked-ruling-pencil}の
左右評価階数の証明）。両議論のこの部分は相手の四次元切断との一致を使わない。
後者では元の点配置の直線切断長が3以下なので次数も3以下である。

混合成分の一般点の全二十印と全位置で三次曲面が滑らかなことは
定理\ref{thm:marked-cubic-surface-pairs}による。従って自己Extは次数0,3のみ各一次元、
Serreによる3移動と合わせて球面性が従う。異点間のExtは全て零である。
点加群の直和の表現核は各隅ごとに点列評価の核である。二次まで単射であり、
三次も系\ref{cor:finite-full-triple-regular-presentation}により一般には全位置で単射である。
\end{proof}
